\section[第三章\q 基本分析]{第三章\\ 基本分析}

\subsection{价格、价值与价值观}

{\kaishu
世界因观察者的不同而不同。

\hfill ——金融帝国

股市的上涨与其说是经济的好转，不如说是交易者更乐意于看到好的一面。

\hfill ——金融帝国
}

所谓价值规律，简单说就是：价格围绕价值上下波动。价格很简单，就是我们电脑屏幕上跳动的数字。那么价值究竟是什么呢？按照某国外书籍的说法，所有的买家都是认为价格低于价值的人，而所有的卖家都是认为价格高于价值的人。那么价值必将是由众多因素决定的，同时这些因素又因观察者的角度不同而不同，并且决不可能存在价值等于什么的绝对公式。由此看来，市场存在的条件似乎就决定了参与者对于价值的“打分”，从根本上绝对不能是相同的！等一下，如果价值也是需要观察者来打分的话，那么这时价格与价值的界限已经开始变得模糊起来。

《资本论》忽略了资源的价值、科技的价值、资本的价值，而完全用社会必要劳动时间来解释价值。也许资本的祖先确实是社会必要劳动时间，但黄金、石油等资源的价值又如何解释？爱迪生数量庞大的发明又价值几何？在我看来所谓价值，不过是为了整体的理论体系而主观添加的一个概念，并且认为价值量应该是绝对的（不受观察者参照系的改变而改变）。也许后来意识到价值不可能存在什么绝对的标准，所以才引伸出所谓价值观的问题。至少在其看来，价值观是有好坏之分的，并且存在一个特殊（正确）的价值观。这说明买家和卖家达成一笔交易时，他们其中必然有一个人的价值观是错的。

这不由让我想起经典物理学，为了解释光的传播而引申出一个绝对静止（不受观察者参照系的改变而改变）的“以太”。我们可以清晰地发现，从亚里士多德的哲学，到欧几里德数学，再到牛顿的经典物理学，最后到古典经济学，从本质上来说是一脉相承的，共同组建了当时一种和谐的理论世界。他们都认为这个世界是简单的、绝对的、刚性的，并且不因任何观察者的不同而改变。

上个世纪一位叫做爱因斯坦的伟人松动了经典物理学大厦的一块砖，最终导致了整个大厦的崩塌。遗憾的是，在哲学观和世界观上，当今人们脑子里还更多地保留了亚里士多德的遗风。爱因斯坦告诉我们，相对性原理；量子理论告诉我们，测不准原理。其中无不都把观察者摆到一个最为核心的位置，世界因观察者的不同而不同！

对于无数的观察角度而言，他们应该都是等效的，而不会有任何一个角度是特殊的。简单说，世界在不同观察者眼中是不同的，而所有观察者的观察并没有优劣、对错之分。

如果把相对论的思维方式引入经济学领域会出现这样的结论：价值对于不同环境、理念的人来说可以是不同的，而其中的任何一个结论都没有对错之分。在相对论中唯独光速是绝对的，不受观察者的改变而改变。那么在经济领域只有价格是绝对的，同样不受观察者的改变而改变。这样就能推导出一个有趣的结论：价格是绝对的，而价值是相对的；价格是真实的，而价值是虚拟的。并且价格是唯一的，而价值可以是无限的，同时这无数个不同的价值中并没有哪一个是特殊的。交易者对市场的感受不过是：自身价值观形成的相对价值与市场绝对真实的价格之间的背离程度。

价值观是抽象的，它来源于每个人头脑中的信念。有意思的是，我们每个人头脑中的价值观也不是静态的，相反它是极不稳定并且变化多端的。事实上，能够影响市场的因素有无限多，而所有投资者考虑到的问题却是有限的。那么从“冰山原理”来看，交易者的信念都是断章取义的、都是盲人摸象的。市场是什么并不重要，重要的是交易者乐意于相信什么。对于市场的波动而言，与其说是因为价值的变化，不如说是价值观的变化。

投资者的价值观无疑在受价格波动的影响而改变。对于大多数交易者而言，他们都相信市场是合理的信念。如果股市大幅度上涨，那么整个行业的所有角色都会寻找股市上涨的理由，从而忽略股市下跌的理由。在这个有上千只股票的股市上，必然会存在250只业绩增长的股票，同时也必然会存在250只业绩滑坡的股票。如果每天分析、探讨一只业绩大幅增长的股票，那么市场将是一片光明而令人兴奋的；但如果每天分析、探讨一只业绩大幅下滑的股票，那么市场将是一片黑暗且令人沮丧的。

国外某心理学家做过这样一个试验：让被测者先在轮盘上转动出一个数字，然后再回答“非洲有多少个国家”的问题。最后的实验结果是：转动出大数字的被测者，通常认为非洲的国家数量更多。如果人们对一个数字应该是多少没有任何概念时，那么他们甚至会以（明知的）随机数字作为参照物，思考是比这个随机数更大或者是更小。这说明什么？任何交易者在对市场价值进行“打分”时，都不可能不受到真实价格的影响。比如，2005年交易者思考市场的真实价值比998点更高还是更低；2007年又思考市场的真实价值比6124点更高还是更低。按照信心理论的说法，交易者对市场的信心随着上涨而不断增加，随着下跌而不断减弱。这再一次证明，行情的波动是可以影响交易者价值观的！当然交易者价值观的变化，也同样会反作用于价格。

这不由让我想起一个有趣的故事。从前在一个小镇上，只有一个钟表店和一座寺庙的钟声能够确定时间。每天敲钟人都会去钟表店来确定时间，然后敲钟报时；而钟表店的老板每天都根据钟声来调整钟表的时间。如果说，钟表上的时间能够左右钟声，而钟声又能够左右钟表上的时间，那么小镇上未来会出现的时间，绝对是不可预知的。

有意思的是，所谓价值是一个非常广义的概念。他不仅包括能够被经济理论认可的内在价值，而且还包括一些“上不了台面”的价值。比如，炒作价值。1999年小盘基金被炒作到净资产的数倍；2007年的认沽权证再度上演“废纸也疯狂”的一幕；2000年转配股上市被当做一种题材炒作……有时我在想，什么是真理？任何事情被多次重复后都能成为真理！在大牛市中，配股被当做一种可以低价购买更多廉价股票的题材而被炒作；在大熊市中，配股被当做一种上市公司“圈钱”的手段而使投资者避之不及。配股这种现象的好坏并不重要，重要的是投资者会用什么样的思维来理解。由此看来，坐庄的成败并不取决于对价格的操控能力，而是取决于对市场参与者价值观的影响能力。

这时的价值规律将变为：价格是所有参与者价值观的综合反映，绝对正确的价值观并不存在。价格作用于所有参与者的价值观，而所有参与者的价值观又会反作用于价格。所谓价值不过是每个参与者根据自身对市场不完全的认知而在头脑中形成的观点罢了。

难道价值观真的没有优劣之分吗？毕竟巴菲特通过其独特的价值观创造了庞大的利润。如果价值观绝对没有优劣之分，那么基本分析就没有任何意义；如果价值观存在绝对正确的标准，那么基本分析就能必然地预知未来。有时我在想，是巴菲特的价值观最终战胜了资本市场，还是资本市场的现实迎合了巴菲特的价值观。为什么巴菲特出现在美国，而不是经济泡沫破灭后的日本？假设在大牛市中，某交易者使用一种“独特”的交易策略，每当价格从高点下跌10％就买入，结果震荡中的牛市总是迎合这个交易者的交易策略，而使其获利丰厚。这样一来，与其说是这个交易者的交易策略战胜了市场，不如说是市场迎合了这个交易者的交易策略。也许在巴西、墨西哥、日本，都存在同巴菲特做着同样事情的交易者，可他们不但没有脱颖而出，反而亏损累累。

如果说巴菲特能够必然地战胜市场，我觉得有些牵强；但如果说巴菲特的成绩完全来自于偶然，我同样觉得牵强。毕竟巴菲特获取的庞大利润是真实存在的。那么我们妥协一下，没有绝对正确的价值观；但可能存在具有优势的价值观。而这种具有优势的价值观也绝对不是一成不变的，甚至也不是事先可知的。

\subsection{价值投资不等于“占便宜”}

{\kaishu
价值投资不是买入好的股票，甚至不是买入将会变好的股票，而是应该买入比大众预期更好的股票。

\hfill ——金融帝国

如果你无法获得大众无法得知的信息，同时又没有与众不同的价值观，那么基本分析对于你来说将永远是一场自欺欺人的骗局。

\hfill ——金融帝国

单凭市盈率高低，无法判断到底应该买入或卖出股票。在多数情况下，一档股票不会因市盈率很低，即大幅涨升，不要只因为市盈率看起来很划算，即买入一档股票。一档股票之所以看起来便宜，时常有它的道理。

\hfill ——奥尼尔
}

上一节我们探讨了价值的问题，那么我们由此想到：所谓价值投资比拼的并不是股票的价值，而是参与者的价值观。前面我已经说了，价值只存在于每个人的头脑里，而且可以是因人而异的。但这与我们通常所理解的价值似乎不同，在市场上似乎还存在一种“市场公认”的价值。往往“市场公认”价值，都是与价格呈现背离状态的。

如果说影响市场的因素包括已知因素和未知因素的话，那么对于大多数参与者而言，这些已知和未知的因素往往在很大程度上是相同的。对于我们这些弱势群体来说，我们并没有资格亲眼见到我们投资的上市公司，更不用说详细了解公司内部的运作状况。这样一来我们对上市公司的认识几乎完全停留在媒体传播的信息与评论上。在这个资信相当发达的21世纪，铺天盖地的媒体已经将这种“市场公认”价值的集中认同度大大弱化。要知道在20世纪90年代，关于市场的媒体并非很多，那时“市场公认”价值的集中度比现在还要大得多。

通常“市场公认”价值包括股票的业绩报告、主流市场评论观点，等等。似乎这时每股收益等基本面数据已经成了不受观察者角度不同而不同的绝对数据。而市场权威研发机构之间的观点也通常是相互影响的，并且融合成一种市场主流预期。

很多交易者就在这里陷入了一个严重的误区。当市场价格远低于“市场公认”价值的时候，那种符合本能的“占便宜”的念头就油然而生。更可笑的是，还美其名曰价值投资。他们一厢情愿地认为，基本分析大师也是在做同样的事情。那么，请看看历史上最为有名的基本分析大师的投资案例吧。

彼得·林奇是美国家喻户晓的基本分析大师。他分散投资的风格独树一帜，他在职业生涯之中，轮番持有过1400多只股票，远远超过同一时期的巴菲特。20世纪80年代，他将旗下麦哲伦基金5％（这是证监会所允许的最大持股比例）的资产买入一只股票。按照林奇的说法，如果法律允许的话，他会把基金资产的10％到20％全投入进去。大家知道林奇罕见的重仓行为买入的是一只什么样的股票吗？那是一个持续低迷的汽车行业中，一只濒临破产的股票——克莱斯勒！最终管理巨子卡特让克莱斯勒起死回生，而林奇先生在克莱斯勒上投资获得了50多倍的投资收益。

20世纪60年代初，巴菲特所作的一个重大决策，便是大力投资美国运通公司。在该公司爆出“安东尼·迪·安杰里斯色拉油丑闻”时，产生了数十亿美元损失的责任。这次危机的爆发有可能卷走全部的股东权益，使公司的净资产变成负值。而巴菲特将巴菲特合伙公司40％的净资产、价值约1300万美元投资于美国运通公司，买下了美国运通公司5％的股份。如此一来，巴菲特打破了他在一次性投资中使用其合伙公司资金不超过25％这一原则，据说巴菲特合伙公司卖掉这些股票净赚2000万美元的利润。我记得，巴菲特曾经戏称：“如果有什么好消息不必对我讲，但如果有什么坏消息请一定要告诉我。”

最后，我给大家讲述一个自己亲身经历的故事吧。2001年时，我还没有放弃对确定性的追求，希望自己所做的每一次交易都是正确的。在大盘两次跌破2000点，我经历了两次止损无效后，我开始决定进行长期投资。那时我买入的是一个经过深幅下跌，并且其后又进行了长达两年筑底过程的股票。那就是有很高每股收益、未分配利润、净资产的股票——600788达尔曼。我没有保护性止损，而我的持股很快从15元跌到了8元多。我记得那时该股的净资产好像是六七块钱，每股收益是0.78元。我当时的一个念头挽救了我的投机生涯。我想无论如何，这么好的股票都不应该跌到如此低的价格。那么就只剩下一种可能——业绩虚假！我在将近亏损了50％以后，出清了所有的持股。此后，该股仍旧快速下跌，并且跌破净资产。那时我更加深信该股一定存在业绩虚假的问题！果然，某日该股爆出巨大的丑闻，奇高的净资产一下子变成负数！很快，该股最终跌到三板市场的0.2元！

我想告诉大家，一只股票正在下跌并不可怕，但如果市场上没有人知道该股为什么下跌才是最可怕的事情。无论是巴菲特还是林奇，他们买入的股票是经过大幅下跌的股票。但请注意，市场上几乎没有人不知道这些股票为什么而下跌！几年后我听到一个故事，一位老大爷在达尔曼跌到6元附近时，用全部几十万的资金买入了这只业绩优良并且跌破净资产的股票。这让我想起经常在马路上出现的一幕，一些外地人拿着一部崭新的价值2000元的手机，以200元的价格向路人出售。不知道你遇到这样的事情后会作何感想呢？是认为这种事情不符合常理而肯定是一个骗局，还是认为这正是占这个傻子便宜的好时候呢？

请记住！价值投资根本不是一种投资理念，而是一种投资能力。基本分析大师能够大获全胜并非因为价值的优势，而是因为价值观的优势。我的意思是说，基本分析大师必然是那些价值观不同于“市场公认”价值的交易者。如果你的价值观与“市场公认”价值并没有任何不同的话，那么千万不要自欺欺人地去进行什么所谓的价值投资。

分析师对市场的观点我们很难找到长期可靠的资料，所以就很难作历史性统计。但如果找一些业绩曾经大幅波动的股票，并且把业绩和价格做成走势图叠加来看，你会发现不是股价跟着业绩走，而是业绩跟着股价走。比如600886，1997年每股收益1元，而后成为连续亏损的ST股票，最终通过重组又一飞冲天。请你看看这十年来是否是业绩跟着股价走呢？如果你还不明白其中的奥秘，请再看一次“冰山原理”。

\subsection{基本分析与自欺欺人}

{\kaishu
对于大多数的股民而言，只有亏损的股票才会进行长期投资。

\hfill ——金融帝国

如果你找到的所有原因都能指向同一个结论的话，那么这通常说明你已经失去了客观。

\hfill ——金融帝国
}

价格是一种买卖双方达成的共识，也就是说买家和卖家通常都有支撑自己行为的理由。除非价格等于0，否则你总能找到一些看多和看空的理由。理论上讲，如果你不进行人为筛选的话，那么你找到的理由绝对不可能都指向同一个方向。如果一位交易者能够真正客观地搜集可能影响未来价格的潜在因素的话，那么当你找到10种原因的时候，通常应该是原因1、2、3、4、5指向未来上涨；而原因6、7、8、9、10指向未来下跌。这时你的综合分析基本上是没有意义的，因为你已经陷入了无谓的矛盾中。当然也许你找到的全部原因不可能绝对多空相等，但同样也绝不可能是指向同一方向。否则，你很可能是在为了迎合自己所处的位置，或者是心理需求而刻意对原因进行了筛选。

在我看来，一位出色的基本分析师应该具有这样的能力：在同一时点上，写出两篇有理有据、逻辑清醒、极具充分说服力的文章。而这两篇不同的文章，却是一篇看多、一篇看空的。很多网友在看过我的文章后，会感到不知所云地质问我：“你到底想说明什么？你到底是看空还是看多？”或许，对于更多的读者来说，没有明确观点的文章绝对不是什么好文章。我曾经听过这样的一个故事，一位科学家来到北极圈内，多年来一直观察那里的生物。有人问他，你研究了这么多年，到底想证明什么呢？而他回答道，我只是在单纯地观察，并不想证明什么。很显然，更多的人并不习惯于通过不断观察而总结出一些规律，而是希望通过观察来证实自己先入为主的结论。不夸张地说，即使是我对市场的观点，也绝无可能不受到持仓位置的影响，所以不得不通过止损点来保护账户中的资金。从某种意义上来讲，刻意远离市场信息，对交易而言是很有意义的。我始终认为交易是一个人的事情，但没有人能够完全避免他人对自己的干扰。不信的话，咱们做一个实验：观察你的呼吸，并且不改变原有的呼吸节奏。很显然，这是不可能做到的。观察行为会使你的呼吸节奏变快，或者变慢。因为我有非常特殊的交易环境，从而对这方面深有体会。刚进入股市时，对于只有十几岁的我来说，完全避免爸妈的干扰是一件不可能的事情。对于爸妈的建议而言，我很难做到不受其影响。当然倒不是说，我会听取爸妈的意见，而是必定会落入“顺从”与“逆反”的两个极端。尤其是当我“逆反”地坚持自己的思路，但市场的运行却完全符合爸妈的建议时，很可能会发生更加微妙的变化。这时我或许变得更加“逆反”，或者彻底“屈服”。无论如何，我都不可能做到把建议当做并不存在！

在我看来，善于进行基本分析的交易者，通常都是敢于为单一理由而下注的人。某香港人看到中国内地报纸关于越南问题的报道中出现“是可忍，孰不可忍”的标题后，预期中越战争的爆发而大量买入黄金；巴菲特考虑到运通仍旧垄断全国旅行支票业务，而大量吃进该股……道理很简单，也许你确实能够找到一个被市场暂时忽略的重要因素，但如果说你找到了N个被市场所忽略的重要因素就显得有些天方夜谭了！

1997年我刚进入股市的时候，买入了一只叫做大庆联谊的股票。可惜的是，自从我买入后该股就不断下跌。当时我不肯接受亏损的现实，而是更乐意相信该股未来会大幅上涨。所以我就拼命寻找该股看好的理由，甚至还拼命寻找大盘看好的理由。现在想想很可笑，似乎整个股市都应该为了迎合我的需求而运行。随着我的持股价格一天一天下跌，我就越是拼命寻找自欺欺人的理由。从而用“好的股票是可以进行长期投资的”等名言来反复安慰自己。那时无论何时何地，只要有人问我如何看股市，我肯定说看涨！有点意思，不是因为我看涨才买入股票，而是因为我买入股票后才只能看涨。

如果有哪位还没有进入市场的读者看到我这篇文章的话，肯定会笑我傻得可爱。但我相信真正成熟起来的股民，在看到这篇文章的时候应该能够回忆起什么。我不想总是说交易有多么困难，但事实上，在市场中做到“客观”二字都不是几年工夫就能做到的。又有多少股民的观点，能够真的绝对不会受到现在自己位置的影响呢？

人嘛，就是这样，只会去相信他乐意相信的“事实”。其实我们每个人随着年龄的增长，都会慢慢形成自己的信仰体系。而人年龄越大，这种信仰体系就越难改变。这是因为我们复杂的心灵已经慢慢形成一种对固有信仰的保护机制，他会把有悖于自身固有信仰的信息抵挡在外。当你和一个中年人交谈时，如果你说的话使他内心表示认同，那么他会表现出一种轻松愉快的表情；如果你说的话使他的内心不认同，那么他也许会和你争论，或者表现出轻蔑一笑；而如果你说的话，能够让他的内心认同，但又有悖于他的信仰体系时，他会启动自身的保护机制，从而表现出一种视而不见的烦躁。如果你越是把明摆着的事实展现给有眼睛就能看得到的他，那么他就会更加表现出巨大的烦躁与不安。这种心理现象，反映在交易者接触到自己不愿意接受的现实时，同样表现得淋漓尽致。

如果这种自欺欺人已经成为一种习惯的话，那么就相当于是在对自己进行千百次的心理暗示。最终交易者就会把他乐意相信的“事实”，当做理所应当的真相。天哪，交易者自己把自己给欺骗了，而且被欺骗得深信不疑。最终，当现实与信念达到一种南辕北辙的程度时，交易者早晚都会有信念崩溃的那一天。

交易者还有一种倾向，就是希望掌握所有的分析方法。但可笑的是，所有的分析方法在他看来，都只不过是让自己随心所欲的交易行为表现出合理化的工具。人性通常会使我们更乐意于兑现获利的股票，而持有亏损的股票。如果某位投资者买入后立刻就出现利润的话，那么他会把“落袋为安”“盈利的交易是不会让交易者破产的”“做交易最忌讳贪婪”“聚沙成塔”“到达阻力位”“指标超买”等一系列交易“名言”来做为他卖出股票的理由。也许这时交易者获利兑现并不需要找任何理由，反而获利就是能够证明自己能力的最好案例。但如果同一位交易者买入后就立即遭遇大幅下跌的话，那么他会把“价值投资”搬出来当做继续持仓的借口。天哪，什么叫价值投资？按照巴菲特的说法，好公司的股票应该是永远不该兑现的！盈利了就兑现，亏损了就耍赖的交易“策略”，绝对不是什么价值投资。我感觉巴菲特之所以成功，应该能够归纳为两种能力：一种能力是选股的能力，另一种能力是持股的能力。事实上，我觉得后一种能力更为重要。请注意，是持有获利股票还能坐着不动的能力！有兴趣的读者朋友可以作一个这样的统计，验证一下身边的股民持股时间与盈亏状态的关系。看看获利的股票最长持有多长时间，而亏损的股票又最长持有多长时间？事实上，大多数股民之所以亏损，输就输在这里！交易其实很简单：像对待亏损那样对待你的利润吧！

\subsection{基本分析更适合谁}

{\kaishu
对于分析方法而言，并没有所谓的优劣之分，重要的是这种方法是否适合你？

\hfill ——金融帝国

投资想赚大钱，必须有耐性。换句话说，预测股价会达到什么水准，往往比预测多久才会到达那种水平容易。

\hfill ——费雪
}

在我看来，所有交易方法都不存在绝对的优劣之分。如果我们把交易方法看做砖块，那么你建立起来的系统大厦的好坏并不取决于砖块，而是取决于砖块的组合。因此，本节我们就来看看什么样的交易者更适合使用基本分析。

从国外的基金业来看，资金量越大的交易者越重视基本分析；操作周期越长的交易者越重视基本分析；基本分析大师通常产生于股市而不是期货市场。从某种意义上来讲，资金量的大小与操作周期的长短也是相辅相成的。我们很难想象用几百个亿的资金，在市场中玩抢帽子的游戏。越大的资金量，在市场中翻身所需要的波动幅度也就越大，并且市场波动幅度又同持有时间成正比的。在期货市场由于保证金制度和合约有效期的问题，使得基本分析大师很难施展拳脚。甚至从另一个角度讲，期货市场中的强势群体更难获得绝对的信息优势。

很多交易者都误以为资金量越大越容易获利，其实不然。当资金量达到一定程度以后，技术分析基本上就不再有用武之地。靠止损来防范风险的策略，基本上也没有任何的可行性。某国外交易书籍描写道，当一位基金经理看着持股暴跌的时候，除了祈祷外不可能去做任何事情，毕竟上百亿的抛盘在短时间内止损是难以想象的事情。这也是基金经理需要用股指期货来实现套期保值的原因。在我看来，一个亿以下的小资金靠趋势跟踪就能获得不逊色于巴菲特的战绩。我个人可以靠“试错”来完成建仓行为，可对于大资金而言这同样是难以想象的。对于资金管理者来说，资金量越大，通常收益率就越低，但同时稳定性就越好。事实上，大资金也有羡慕小资金的地方。

当然在中国还有一类人会非常重视基本分析，那就是市场操纵者。他们分析基本面只不过是为了找到一个富有想象力、便于炒作的概念。他们并不指望这种概念最终能够对业绩产生多大的影响，而只是为了找到一种能够让散户认同的借口。毕竟，只有散户的大量认同，庄家才能成功地实现撤庄。

基本分析的进入门槛是非常高的。据说在一个大型券商的市场研发部门，对于一个行业就要组建一个分析小组。这种专业性的调研工作是我们中小投资者所难以效仿的。这要求具备多年积累下来的专业知识；需要从信息采集到信息分析各个环节的协同作战；甚至还需要高昂的调研成本。以下是关于林奇的一段描述：林奇是一个工作狂，每天的工作时间长达12个小时，对其所做的一切显示出一种着了魔的狂热。也许在投资界没有人比他工作更努力，也没有人比他阅览所及的范围更广。他每天要阅读几英尺厚的文件，他每年要旅行16万公里去各地进行实地考察，此外，每年他还要与500多家公司的经理进行交谈，在不进行阅读和访问时，他则会用几小时、几十小时打电话。也许只有庞大的资金，才能分摊巨大的调研费用。而我们这些业余投资者，既承担不起庞大的研发费用，也不可能付出如此大的精力。

股市上的风险分为系统性风险和非系统性风险。从中短期来看，股价的走势和大盘具有高度的相关性。这说明只有操作周期大到一定程度以后，非系统性因素才能大于系统性因素。而非系统性因素也不光只包括公司的基本面因素，同时还包括市场的潮流、庄家的炒作行为等。同样只有操作周期大到一定程度以后，基本面因素才能大于人为因素。

道氏理论指出，趋势分为长中短三个层次。长期趋势必然是由基本面的因素来决定的；中期趋势中人为因素只能干扰趋势；而短期趋势是完全可以人为操纵的。这同样说明基本面因素只有随着时间周期的增大才能占据绝对的主导地位，而其他因素会随着时间周期的增大而衰减。这就好像万有引力一样，它是一种非常弱的力量。在微观世界中万有引力几乎可以忽略不计，而强力和电磁力才起决定性作用力。但从微观世界转换到宏观世界的话，万有引力才是维系整个宇宙最重要的一种力，而强力和电磁力随着距离的增长，而快速衰减。我觉得这种比喻很恰当。很多中小投资者其实根本就是生活在微观世界里，对于微观世界的现象他们不考虑强力、不考虑电磁力，而恰恰乐此不疲地考虑万有引力。

那么多长的操作周期，基本分析的力量才能够成为市场上的主导力量呢？在我看来，最低限度也必须包容市场走势与基本面呈现非相关的时间跨度。中国股市在1997年至2001年的四年间，基本面几乎与走势呈现负相关的特征。通常来说，交易方法必须要包容成本。那么靠基本分析实现获利的交易方法，至少应该需要五年的时间跨度。

事实上，基本分析最难解决的就是时间问题。也许预测股市要发生什么事情，要比预测什么时候发生简单得多。我回忆了一下1998年看过的很多书籍中提到的股票，有一本书的作者通过基本面分析非常看好山西汾酒，说他准备长期投资，并且用每月工资中的一部分钱去加仓。2006年山西汾酒确实成为了黑马，但愿他能坚持这8年的时间。还有一本书的作者看好盐湖钾肥，指出该公司拥有全国95％的钾肥资源和产销量，所以成长性极佳。事实上到了这几年该股确实异常强劲。1999年我曾经密切关注过长安汽车，该公司的质量意识非常强，具有大企业的风范，并且汽车进入家庭也是大势所趋，可惜该股到了2003年才一飞冲天。当然这都是我找出的一些成功的案例，1998年还有很多基本分析者看好湘酒鬼，遗憾的是，该股最后业绩滑坡到ST板块，2005年价位已经跌破3元。

在中小投资者通常的操作周期中，坏股票可以让你获利丰厚，而好股票也可以让你输掉裤子。在我看来，交易者应该拿自己的长处同别人的短处来较量。否则，非要不顾自身环境，而一味追求基本分析大师的交易思路，往往会得不偿失。世界上没有最好的交易方法，只有最适合自己的交易方法。在我看来，基本分析实在是一场不平等的游戏，但在价格面前绝对是人人平等。我始终想不明白，交易者为什么对这种不平等的游戏总是乐此不疲呢？作为一个连真相都未必能够知道的你，又如何能够打败可以改变真相的对手呢？我与其相信巴菲特是一位持股傻等的世外高人，还不如相信巴菲特是一位出色的资产重组高手。有一句关于感情问题的名言：没有人能够伤害到你，除非是你自己想受到伤害。同理，没有任何一种交易方法能够伤害到你，除非是你想受到伤害。如果你不相信基本分析的话，那么基本面的骗局无论如何都不可能伤害到你。当然，技术分析也有骗局。但我能够接受、包容这种缺陷，我可以心甘情愿地反复做着无效止损，而不会去责怪任何人。如果你心甘情愿做一位基本分析的信徒的话，那么请你一定要包容基本分析的缺陷，当受到伤害的时候请一定不要抱怨。我经常受到趋势跟踪的伤害，但我从来都不抱怨。我最瞧不起一边指责股评，一边去听股评的交易者。可能他们一辈子也无法懂得“取舍”二字的含义；一辈子也无法懂得“责任”二字的含义。如果中小投资者非要运用基本分析来进行投资，那么至少你应该答应我，绝对不要将获利的股票在5年内兑现出局。

在本节的最后，我继续给大家讲一个故事吧。1999年底的中期熊市中，市场上一只叫做ST渤化的股票不可思议地走出了连续涨停的行情。按照交易所的规定，连续三个交易日达到涨跌幅限制的股票，上市公司必须公布“澄清公告”。在连续十个涨停板的过程中，上市公司连续三次公布“公司没有任何应披露而未披露的信息，请广大投资者注意投资风险”。直到最后一个涨停板时，上市公司公布了“资产重组方案仍在洽谈中”的信息。更有意思的是，此后该股走出了连续五个跌停板的走势。与此类似的案例还有很多很多，阿城钢铁资产重组为科利华时，也是在股价翻番后，中小投资者才得知重组的真相。善良的中小投资者们，请相信我：在基本分析面前，我们毫无优势可言。

\subsection{基本分析与技术分析的兼容性}

{\kaishu
在股市上，当你有了绝对的止损点后，就不用再考虑其他防范风险的方法了。事实上，止损点换来的绝对不只是对风险的控制，他对获取更大的利润也有非常重要的意义。

\hfill ——金融帝国

我不知道天有多高（利润），但我知道地有多厚（由止损决定的亏损）。

\hfill ——金融帝国
}

自古以来，交易者就分成了基本分析派与技术分析派两大阵营，从而展开了一场孰优孰劣的争论。有人提出，为什么就不能将基本分析与技术分析兼容起来考虑问题呢？难道非要在技术分析与基本分析中进行二选一吗？

基本分析从价格运行的根本原因入手，在运行方向与幅度的判断上有一定的优势；而技术分析从价格波动的本身入手，对于具体点位的把握有一定的优势。那么，如果用基本分析进行品种和交易方向的选择，然后再通过技术分析来寻找具体的买入点位岂不是更好？理论上讲，这样的方法是可行的。当一只股票基本面显示该股应该大幅上涨，而现实的股价却在不断下跌时，我们不妨先不要逆势买入，而是等到该股技术上摆脱下降趋势再买也不迟。要知道我们手上那点可怜的资金，完全没有必要像基金管理人那样逆势吸纳。

至此没有任何的问题，但是如果买入后该股又继续下跌该如何呢？基本分析的本质就注定了其不可能存在弹性，我们很难将止损思维与基本分析兼容起来。基本分析是因为“便宜”而买入，那么我们又如何能够因为更便宜而卖出呢？按照巴菲特的说法，当其买入的股票出现下跌时，不过是给他一次以更便宜的价格买入的机会。由此看来，基本分析大师通常是不止损的。

那么，基本分析大师又是如何防范风险的呢？按照林奇的说法：假如你手中的十只股票有三只是大赢家，他们就能弥补你其中一两只的损失和六七只表现平平的股票。对于巴菲特而言，也有一次性投资中使用其合伙公司资金不超过25％的原则。由此看来，基本分析大师是靠分散投资来防范风险的。如果持股时间超过一定尺度的话，那么系统风险基本上是不存在的，毕竟指数没有不能逾越的顶部。那么基本分析者，完全可以通过分散投资来规避非系统风险。

在股市上，如果你有绝对的止损点，那么分散投资的意义不大。但如果是在期货市场上，即使你有止损点，也必须要考虑到分散投资的问题。事实上，方法并不是最重要的，最重要的是如何搭配，并且让这种搭配能够适合你的自身要求以及市场环境。对于交易而言，你必须将买入、持仓、卖出组成一个环。即使这样也不存在能够保证一次见效的交易策略，那么我们还要将这些众多的环组成一条链。一个人交易做得好不好，关键是看他的交易链是否很好。

既然没有能够为单次交易负责的方法，那么技术分析只能靠多次运用来实现它的价值。而多次使用技术分析，就等于形成了一个交易系统。这时就出现了盈亏同源的问题，利润不再是利润，而只不过是收入；亏损也不再是亏损，也只不过是成本。这时如果我们在付出多次成本后，又放走了属于我们的利润，那时我们最终只能是出现总体亏损。我们在规避了亏损的同时，可能也会放走潜在的利润；而如果放走了利润，就绝对不是没有赚钱的问题。因为我们如果不能用必然的利润来对冲必然的亏损，那么我们只能出现亏损的结局。对于交易系统而言，并不可能每次买入信号都能够有基本面的配合。

即使是最好的基本分析者，也不可能对所有品种都能够做到清晰的分析。这时采用基本分析就会存在放走黑马的可能。如果你只使用基本分析的话，那么放走其他的黑马并无可厚非。但如果你放走的黑马，恰恰是你由技术分析构建起来的交易系统所能够捕捉的话，那么这对你而言就是一个非常重要的损失。按照“盈亏同源”的理念来看，放走利润和出现亏损并没有什么本质的不同。要知道，放走一次利润可能会出现数次亏损无法对冲。巴菲特可以放走网络股的利润，毕竟他有属于自己的盈利模式。但趋势跟踪者如果放走网络股的利润，那么他该如何对冲曾经多次无效止损的成本呢？

由此我们可以得知，将基本分析与技术分析融合起来应对某次特定交易的可行性是无可厚非的，但如果将基本分析与技术分析结合起来组建一个能够循环使用的交易系统，就绝对不是那么简单的事情了。何况每一种交易方法不但有其优势，同样也有其成本。如果我们不伦不类地将不相干的交易思路融合至一起，就会出现众多基本面因素加上一个紧密止损点的可笑局面。要知道任何基本分析都不能保证价格不会反向波动百分之几。

在我看来，每一种特定的交易方法都应该有其适合自身的风险防范策略。风险防范策略从本质上来看是一种成本。以止损策略为例，这种思维不过是把可能的错误当做错误来处理，自然这里就会存在“误杀”的可能。风险防范策略本来就不是越多越好，而是只要能够对你的整体系统有效就可以了。比如，在雨天出行的时候。我们可以选择直接打车；也可以选择打一把雨伞来乘坐公交车；还可以穿上雨衣骑自行车。我们完全没有必要穿上雨衣，再打一把雨伞，然后打车出行。

我相信对于任何交易者而言，最难接受的风险防范策略就是止损点，这意味着交易者必须放弃本能对确定性的追求。我也是经历了数年的寻找用其他方法来防范风险都无效之后，才心甘情愿地使用“无条件”的止损策略。最早我用买入业绩好的股票来防范风险，结果业绩好的股票也可能出现大幅下跌，这时我无法防范业绩滑坡和虚假的风险；后来我采用买入长期下跌的股票来防范风险，结果弱的股票通常会更弱；最后我用择强汰弱的方法来买入股票，结果同样还是不能够保证不会大幅亏损。没有办法，我只能放弃确定性而接受无条件止损。要知道止损策略的成本是相当高的，我的经验表明超过半数的止损通常都是无效的，而且超过半数的买入（开仓）都会产生止损的结局，甚至卖在地板价上也是经常的事情。但与此同时，我也获得了一种得天独厚的优势。业绩差的股票我敢买；价位高的股票我也敢买；大幅上扬后的股票我照样敢买；几乎没有我不敢买的股票。道理很简单，无论一只股票长期的风险有多大，对于我来说短期的风险都是一样的，那就是止损点的幅度。而理论上来讲，除价格走势本身外，没有任何选股策略能够对减少触及止损点的次数有效。事实上，市场上很多快速上扬的股票、高市盈率的股票、价位高的股票，他们的庞大利润经常都是无人认领。一个简单的止损策略，就能让我不用忌讳这些风险。后来我才知道，原来这些股票不但不是我们想象的那样具有巨大的风险，反而是一种最安全的选股策略，至少这样的股票通常更有大幅上涨的潜力。哪一类股票中持股的散户少，那么哪一类股票的风险就小。根据我这十多年的经验，当牛市结束的时候，价位越低、涨幅越小的股票，下跌的速度就越快。道理很简单，在每一轮大牛市中，中小投资者都会集体上演持仓向弱势股集中的闹剧。毕竟，很少有人能够有勇气在更高的价位将抛出的股票买回，而只能不得不寻找更弱的股票。

从交易者的信念体系来看，基本分析与技术分析也是截然不同的。技术分析属于统计学范畴，通过具有概率优势并且富有弹性的方法来实现获利。对于一位好的技术分析者而言，他可以完全不用考虑未来的走势如何，利润来自于不断的试错。而基本分析则不同，他是通过指出未来市场“应该”运行的方向与幅度而实现获利的分析方法。由此看来，基本分析者是必须能够“看懂”未来走势的人。这时出现了一个不容忽视的问题：当你以试错的心态来实现止损时，同以分析和预测的心态来实现止损时，心理的反应是截然不同的！对于分析和预测而言，止损能够给你带来一种挫败感。一次止损还好，但如果是连续多次止损的话，那么这种挫败感将会变得难以让人接受。从我的成长经历来看，让我亏损最惨重的交易，通常都是我最有信心的交易。当我得知越多看好的理由时，我就越难实行无条件止损。毕竟应该止损的时候，我的那些看好理由通常都还存在。如果不等到价位低得已经伤及元气的话，那么这些上好理由通常是不会消失的。人都是有感情的，当你由于过多地了解一只股票的优点而产生感情时，那种保持弹性的转变将会越来越困难。

对于我们这些个人投资者而言，如果你没有与众不同的信息源或价值观的话，那么当基本面与现实走势分道扬镳的时候，永远相信现实走势！市场既是运动员又是裁判员，毕竟所谓的“理”是由市场来定的。当成为废纸的钾肥权证还有人愿意用0.1元的价格来购买时；当认购权证的价格比正股还高时；当豆粕的价格比黄豆还高时，你会作何感想？如果是保证金交易的话，除非你有足够的资金和足够的耐心支持你与“不合理”的波动对抗，否则还是不要和市场讲理。如果市场已经不理智的话，那么我们也没有必要非理智不可。在市场中交易时间越长的人，往往就越能敬畏市场的力量。

\subsection{消息面与未来走势}

{\kaishu
一根火柴能够引发一场森林大火，并不取决于这个火柴，而是取决于当时森林的状况。

\hfill ——交易名言

管理层调控市场的政策，往往与管理层的愿望相同，但又与管理层的预期相反。

\hfill ——金融帝国
}

如果市场上有利好（空）消息公布，那么我们可以预期第二天会出现高（低）开的行情。但是高开以后市场会如何运行，就不是消息本身能够决定的。这说明消息本身并不能直接给我们带来任何操作提示。我们只不过提前知晓第二天可能出现的盈亏，而我们对这种盈亏根本无能为力。

一根火柴能否引发一场森林大火，并不取决于这根火柴，而是取决于当时森林的状况。没错，同样的消息在不同的市场环境公布，产生的结果可能也是截然不同的。在我看来，市场对消息的反应要比消息本身（对交易者的提示）更加重要。下面我们来看看十年来，股市消息面与未来走势特征的分类。

第一类，组合调控，最终见效。案例一，1996年，管理层用“十二道金牌”的组合拳给市场降温；案例二，2005年以前，管理层陆续推出多种利好为持续低迷的市场打气；案例三，2007年，管理层用多种政策调控近乎疯狂的中国股市。

此类消息面都是管理层对市场的调控，并且最终都起到了一定程度的作用。但不可否认的是，之前的多次调控都在反复被市场嘲笑。渐渐地，交易者已经对调控开始麻木，而只有最后一根稻草才压倒了骆驼。记得2005年时，市场流传着这样一种说法：千万别再出利好了，管理层出一次利好，大盘就下一个台阶。这多少有点像“狼来了”的故事，不等到没有人会相信的时候，狼是不会真的到来。由此可见，直接按照消息指示的方向操作往往并不明智。但我们也不能忽略管理层的意图，通常管理层是不达目的不罢休的。有意思的是，无论是1996年底的暴跌还是2007年5·30调整，都没有使得市场最终形成顶部，而只是使得市场出现了快速而猛烈的调整。这或许是因为市场的顶部阶段需要沸腾的市场人气，而管理层的不断调控反而使人性中的疯狂受到压制。

第二类，底部利空。案例一，1997年2月邓小平逝世；案例二，1998年1月香港百富勤清盘；案例三，1999年5月大使馆遭导弹袭击；案例四，1999年12月转配股上市、国有股高价配售；案例五，2002年1月媒体在暂停国有股配售后继续讨论国有股问题；2005年底全流通方案正式推出。

此类消息通常为偶发事件，或者是预期利空最终兑现，极少数为管理层使市场雪上加霜。通常这些消息并非有实质性的内容，但又都被市场过度反应而造成一定程度的恐慌。奇怪的是，超过半数的底部似乎都有利空消息的身影。坦率地说，当时整个市场完全笼罩在悲观的气氛当中，使我多年以后仍旧能够对那种恐慌记忆犹新。由此看来，大幅下跌后当市场对利空消息极其敏感时，直接按照消息指示的方向操作似乎有些显得愚蠢。

第三类，底部利好。历史上真正意义上的底部利好，只有1994年325点时推出的“三大政策”。2005年以前虽然管理层暖风频吹，但市场丝毫都不领情。只有全流通方案出台，使得“第二只鞋子落地”后，股市才开始不断飙升。在我印象中，大盘不断下跌时，管理层推出的利好通常都是一日行情，有时都持续不到收盘。其实道理也很简单，出人意料的利好反而打乱了市场的恐慌气氛，自然就很难形成真正的底部。由此看来，熊市中见利好买入股票同样不可取。

第四类，顶部利空。案例一，1996年底《人民日报》社论；案例二，2001年国有股配售。在我的印象中，好像就有这两次。《人民日报》社论只是让股市快速下跌，但并没有形成最终的顶部。而国有股配售推出之际，市场的反应异常冷淡。有意思的是：6月14日，国有股减持办法出台时，时任财政部部长项怀诚当即表示：“国有股减持是一个利好因素。”2007年5月30日大幅提高印花税的消息，确实让股市出现了多年少见的调整行情。但之前管理层也推出了无数的调控政策，只有这一次才产生了效果。

第五类，顶部利好。案例一，1997年7月香港回归；案例二，1997年9月“十五大”召开；案例三，1998年12月份降息；案例四，1999年9月允许三类企业入市；案例五，2000年以后的《大雨过后是晴天》、以及对股市恢复性上涨的定义；案例六，2002年6月再次重申暂停全流通；案例七，2004年1月底《国九条》的推出。顶部的利好也是很常见的，这与底部利空是对称的，其道理是相同的。

由此看来，消息面的指向同未来中期走势呈现出轻度负相关的特征。一味按照消息指引的方向操作，通常会得不偿失。对消息解读的能力，也不是一两年能够养成的。对于新股民而言，最好的办法就是当消息不存在，坚信走势说明一切。“眼不见、心不乱”，未尝不是明智之举。如果你当消息不存在，那么消息至少不会伤害到你。在我看来，吃过消息亏的股民要数倍、甚至数十倍于占过消息便宜的股民。

当然对于有经验的股民而言，消息最大的价值就在于它是非常不错的“试金石”。用消息来验证市场的强弱，这一点在个股上非常灵验。其核心思想就是：该跌不跌，理应看涨；该涨不涨，理应看跌。以消息面作为参照物，我们就能对比出走势的真正强弱。基本面是静态的，而消息面是动态的。无论如何，他们都适用一种相同的思维方式，那就是“冰山原理”中提出的“A+X=B”公式。已知原因越好（坏），走势越坏（好），就说明未知原因更坏（好），所以后市看空（多）。

其实，我们还可以从另一个角度来看待管理层调控市场的行为。通常来说，管理层调控市场的政策，往往与管理层的愿望相同，但又与管理层的预期相反。简单地说，如果管理层不是认为市场会继续上涨的话，那么又何必出台所谓的利空消息来给市场降温呢？毕竟，管理层不是交易者，其动机并非是获利，而是保证市场的稳定。那么，所谓的调控政策往往都是逆势的，从而使无效调控的次数大大高于有效调控的次数。在上面的案例中，我们只是回顾了市场转折点时的消息面，而更多被镶嵌在趋势中的无效调控我们没有一一列举。反过来说，如果管理层推出的消息与市场的趋势方向相同的话，那么除了说明管理层对于趋势的反转胸有成竹以外，还有其他合理的解释吗？比方说，998点以前管理层为什么不惧怕全流通的正式实施，有使中国股市崩盘的风险呢？

事实上，大部分的调控消息都很难简单地归纳为利好或者利空。比如，国储局抛铜，你觉得是利好还是利空呢？表面上来看应该是利空，毕竟可以增加铜的供应量，对吧？但实际上，国储局抛出的铜越多，就说明库存的铜所剩越少，未来继续抛售铜的可能性越小。客观地说，这样的消息应该属于短期利空、中长期利好。也就是说，如果国储没有魔术棒来变出更多的铜，那么抛与不抛只不过是一种对于短期压力与长期压力的转换。从历史统计中我们不难发现，市场通常将抛储当做利好来理解，而将收储当做利空来理解。事实上，市场上并不可能存在绝对意义上的利好。比方说，“四万亿经济刺激政策”同样具有硬币的两面。否则，这个世界就不会再有什么经济难题了，只要有危机就可以用发钞票来挽救嘛。从某种意义上来讲，救市政策在中短期的确使得股市出现大幅反弹的走势，但也为中长期紧缩的货币政策埋下了伏笔。很显然，管理层也并没有神奇的魔术棒来彻底消除危机，而只不过是按照可接受的方式对危机进行转化罢了。有一种观点认为，不能一味地认为“软着陆”就一定要比“硬着陆”好。日本经济泡沫的破灭应该算是软着陆，而东南亚各国在货币危机中应该是硬着陆。问题是，日本为了避免经济出现硬着陆，而不得不花费超过二十年的时间来消化此前的泡沫。

\subsection{业绩不会永远是市场的主线}

{\kaishu
从宏观上来说，业绩好的股票通常是价格高的股票，但并不等于是未来会上涨的股票。

\hfill ——金融帝国

已有的事，后必再有；已行的事，后必再行。日光之下，并无新事。

\hfill ——《圣经》
}

我承认，对于股票未来的业绩如何，我基本上是无能为力的。所谓以业绩为主线，是指以当时业绩优良的股票作为市场热点。有一点我要说明，“绩优股”与“绩差股”只不过是一种按照业绩好坏的分类。他们同“大盘股”与“小盘股”的分类一样，从宏观上来看对未来走势并没有任何的指引作用。能够左右未来走势的是未来的（超出预期的）业绩。我的意思是说，市场在某一阶段绩优股会表现出众，而另一段时间绩差股又会表现喜人。

我并不爱预测行情，但为了阐述上面的道理，我还是大胆地尝试一下对明年行情的分析（本章节写于2007年5·30调整中。事后来看，预测中的市场现象还是正确的，但时间上就差距很大了。毕竟，5·30调整在很短的时间内就结束了）。我个人感觉，现在的股市行情与1996年的行情有些类似。1996年大盘上涨了两倍，而这两年大盘上涨了三倍；1996年全部股票都出现了一致性上涨，而这两年也是所有的股票都出现了一致性上涨；1996年的末端垃圾股快速暴涨，而今年上半年三线股也大幅飙升；1996年的股市被《人民日报》社论打下来，而今年的股市也被“大幅提高印花税”打压下来；1996年底的下跌中垃圾股跌幅最深，而今年股市的下跌也是三线股首当其冲……

股市在应对“大幅提高印花税”时，三线股夺路而逃，而高价股表现得非常抗跌。这种现象其实就为第二年股市做出了清晰的指引，明年的行情应该是高价股的行情。当然高价股基本上就是绩优股，所以明年的行情应该是以业绩为主线的行情。

回想十年前的市场，我发现一个有意思的现象。1997年以长虹、发展、深科技为首的绩优股大幅上涨，涨幅达到200％以上；与此同时以上海石化为首的低价大盘股都没能创出1996年底的新高；而价格很低的绩差股同样没有怎么上涨，甚至整轮大牛市都没有什么起色。

如果说本轮行情是大牛市主升浪的话，那么其后的五浪通常会有一个板块无法创出三浪的新高。主升浪的特征为全部股票一致性的上涨，所以主升浪比的是持股的耐心；五浪通常是分化性上涨，而绝对不是所有股票都能够有良好的表现机会。五浪通常是一个强者恒强的市场，越高的股票涨得越多，那么五浪比的就是选股的胆量。在五浪中等待补涨的思路，将永远是一个骗局。按照这种思路推导下去，那么分化的五浪谁涨谁不涨似乎已经一目了然了。

难道这真是所谓的“十年一循环”吗？十年前的一幕会在现在重演？中国股市有一个奇怪的现象，价格最高的股票与价格最低的股票比价通常在30倍以内，而从来没有超过50倍。我猜想，明年的行情极有可能大幅刷新这种比价效应。从国际市场来看，这种比价效应要比中国市场大得多。也许未来的绩差股相对于整个市场而言，会表现得惨不忍睹。我做过一个调查：现在很多股民还认为几块钱的三线股便宜，而一百多元的绩优股贵。我相信当绩优股涨到几百元，而三线股还是几块钱的时候，大众的评价就会反过来。

但大家想过这样的问题吗？等到业绩理念被发挥到极致时，等到高、低价股的比价达到三位数时，之后市场会如何呢？看看1998年以后的垃圾股，你也许会明白一些道理。在1998年时，我一直认为只有中国股市“不健康”。业绩好的股票不涨，业绩坏的股票涨，而且市场上经常会爆出财务丑闻。但现在看来，中国股市确实从“量”上有些过分，但本质上与世界股市并没有“质”的不同。美国股市2000年不也是暴炒亏损的网络股吗？美国股市不也出了一个“安然事件”吗？如果本节的假设是对的，那么一年后的市场参与者会认为业绩是永恒的主题、市场不会再有业绩的虚假、甚至好的公司都不会出现业绩的滑坡、暴炒垃圾股的年代一去不复返、现在是一个充满光明的年代……“已有的事，后必再有；已行的事，后必再行。日光之下，并无新事。”《圣经》告诉我们，世界是在一种循环中向前发展。也许你不相信，但我深信时间会证明我是对的：曾经在股市上出现的现象，在一段时间后还会再度出现。即使我不知道还要等待多长时间，但我相信过去的一定还会回来。

本节的预测大家可以当做笑话来听，因为对错根本就不重要。我想表达的是，在微观上交易者可以追随市场的潮流，但在宏观上不要把业绩的好坏与涨跌等同看待，更不要让自己的思维固化。当大家把现在的业绩优良理所应当地等同于未来会上涨的时候，小心市场出现“巨变效应”。当这个市场上股票业绩已经出现巨大悬殊时，赌好的公司会变得更好，不如赌坏的公司会变得不那么坏。《圣经》中说：“一将变成一千；小国将变成强国。”强弱互换是自然法则，“否极”的时候就要敢于赌“泰来”。

\subsection{我眼中的基本分析}

{\kaishu
从某种意义上来讲，在忽略了市场真实价格走势以及参与者思维偏向的情况下，研究基本分析是没有意义的。

\hfill ——金融帝国
}

前面谈了这么多关于基本分析的问题。在我看来，基本分析虽然对于小资金来说意义不大，但却是一种非常有趣的思维游戏。我之所以在测试版后添加本节的内容，是因为我对基本分析的看法有了一些微妙的转变。

就我个人的风格来看，注定无法接受繁琐与复杂的东西，而更喜欢思考一些更接近于本质的东西。或许这并不是巧合，就好像盲人往往都会有很好的听力一样。走了这么多年的系统之路，我发现了一些自己无法逾越的瓶颈。我的计算机知识与数学水平少得可怜，自然就没有办法进行更多的测试，从而寻找市场上更多的微观规律。但与此同时，我却更乐意用与众不同的视角来思考基本分析的问题。当然，我对市场的分析以及预测，并不能直接运用于实战，但我仍旧在享受着思考的乐趣。

如果我们回到2005年那个时点上，并且不受事后合理偏向的影响，那么我们该如何客观定义那时的基本面呢？对于实体经济而言，一个明显的问题就是产能过剩；对于股票市场而言，更加明显的是全流通所带来的巨大股票供给。事后看来，加上超级大盘股的集体上市，股票的供给比那时人们所能够想象的还要大得多。

无论是产能过剩还是全流通，我们都可以简单地认为会造成股票市场的下跌。但如果换一个角度来思考呢？解决产能过剩、实现股票市场的全流通，或许只有靠信贷扩张推高资产价格来解决。简单地说，对全流通问题的解决，必须要依靠大牛市的市场氛围。按照这样的逻辑，就会得出一个有意思的结论：某种经济问题会造成股市的下跌，但解决这种经济问题的唯一方法却是股市上涨。其实，当我们回顾1999年5·19行情时，也会发现这种类似的现象：因为管理层具有迫切的、更多的融资需要，所以必须让股票市场活跃起来。

这是否有点类似于“人择原理”：宇宙为什么是这个样子？基本自然常数为什么不是其他的数值？因为只有宇宙是这个样子，才能够有人类的出现，从而谈论这个问题！因为人类“必须”出现，所以宇宙就只能是这个样子。有点意思吧，原因和结果的位置已经变得非常模糊了！我们用“因果关系”根本无法解释为什么股票供给增大会使股市上涨。但我们却可以用“相关性”来解释：因为股票发行数量与行情的火爆程度具有明显的正相关性。同样的道理，只有物价在飞涨的时候，人们才会买得更疯狂。记得在我小的时候，曾经出现过一次抢购潮，而那时抢购的布料至今还“珍藏”在我家的柜子中。我很喜欢听爸妈讲述那个年代的故事，即使是库存积压的残次品，也会被消费者抢购一空。现在想想看，如果没有通胀引发的抢购潮，那些在我家柜子中的布料又会有什么样的归宿呢？

似乎解决这一切问题的办法就是信贷的扩张，而信贷扩张显然需要依赖获利效应。记得2007年的时候，申购新股的年化收益达到两位数，而我在银行还听说过保年化收益率达到10％的信托产品。要知道，这种无风险收益比银行贷款利率还要高。如果市场参与者是理性的话，那么就应该把房屋抵押给银行来申购新股。最低限度，购房者如果不用尽80％的贷款额度，绝对是一种愚蠢的行为。信贷的扩张可以推高资产价格，从而形成获利效应。获利效应与信贷扩张，可以说是一种“风助火势，火借风威”的正反馈效应。

这似乎就是索罗斯在“反身理论”中所说的“鞋襻效应”。在我看来，索罗斯的反身理论是目前为止，对市场本质描写最好的理论。至少在索罗斯看来，在忽略市场参与者偏向的基础上研究基本分析是没有意义的。甚至，历史进程就是由参与者的偏向所推动的。或许“金融炼金术”的本意在于：意志并不能够使得铅块变成黄金，但市场参与者的意志（偏向）却可以改变真实的世界。

很多人都知道索罗斯质疑市场永远是正确的说法，这似乎触动了技术分析三大假设中的基石。事实上，索罗斯的本意并非是质疑市场超前于实体经济，相反他是支持这种假设的。只不过对于这种现象的解释，他并不认为是市场存在预测的功能，而是因为市场的波动可以改变（左右）实体经济。从另一个角度来说，索罗斯并不像巴菲特那样，在市场犯错的时候去做正确的事。而是在市场犯错的时候“将错就错”，甚至是提前预测市场会犯什么样的错误。这一点在“抵押信托公司”的交易上体现得淋漓尽致。即使他知道这类公司最终会破产，但仍旧不妨碍他先去做多、再去做空。很显然，相对于“真实”的基本面而言，索罗斯更在乎市场参与者会出现什么样的“偏向”。

对于信贷扩张的研究，似乎更应该从市场参与者的认知入手。现在出现了另一些问题：为什么市场会在那个微妙的时点上出现信贷扩张呢？或者说，如果没有出现信贷扩张的话，真实的经济是否会出现另一条发展轨迹？信贷扩张的动力既然是参与者存在偏向的认知，那么这种偏向最终是否真的能够改变历史呢？

对于用“解决产能过剩、实现股票市场的全流通，或许只有靠信贷扩张推高资产价格来解决”的理由来解释信贷扩张的成因，很多人的确有理由不屑一顾。毕竟，更多的人习惯用“因果关系”的逻辑来思考问题。一位朋友曾经这样比喻我：我很像原始人类，将任何无法解释的现象都归为神的力量。我觉得他的比喻很有意思，从而采用这个比喻来阐述自己的观点：无论人类认为地震是因为地壳的摩擦，还是认为地震是神的力量，只要地震是不能够预测的，那么任何的归因（对于行为）都是等效的。即使是现代人在地震发生的当时，除了祈祷还能做些什么呢？即便古人用“风伯、电母”来解释自然现象，也仍旧不妨碍他们在下雨时披上斗笠，这并不需要精通物理知识。如果说人类对于每一项技术的使用，都需要首先了解其中原理的话，那么恐怕人类就无法进化到今天的程度。古人虽然不知道“火”到底是一种什么东西，更不知道“生火”的原理，但仍旧不妨碍通过经验来使用火。即使是对于现代人来说，我们除了知道量子论的精确以外，对其原理仍旧是一无所知。很显然，我们无论是把信贷扩张的起因假设为管理层有意操控，还是神的力量，或者是别的什么原因，对于我们这些市场参与者而言，应该都是完全等效的。

当然，我必须承认上述的分析方法并不是可以预知未来的水晶球。如果我们在2001年因为国有股减持的原因而预测股市会上涨的话，那么就显然是一次错误的预测。这或许说明，即使是“神的力量”，在市场上也并非是无所不能的。很显然，信贷扩张的出现是需要依赖特定土壤的。我始终觉得，大众的储蓄行为才是通胀能够发生的源动力。如果没有上一代人沉淀下来的储蓄，那么恐怕信贷扩张就不会出现，至少是得不偿失的。事实上，在过去的十几年中，管理层为了避免通货紧缩的出现，只能投放更多的货币来抵消民间的储蓄行为。而这种长期投放更多货币的累积效应，再加上居民储蓄回流市场的可能性成为了信贷扩张所需要的特定土壤。由此可见，在市场参与者偏向下形成的价格，并不可能天马行空、不受限制地影响到“价值”。但此后，在正反馈效应下形成的信贷扩张，能够演绎到什么样的程度就更多地取决于参与者的偏向了。有意思的是，参与者的偏向给市场带来的影响，在一定程度上呈现出一种不可逆性。

比方说，参与者的偏向（居民不理性的借贷、银行不理性的放贷）可以使房价从两千元每平方米飙升到两万元每平方米，但管理层绝对无法接受房价从两万下跌到两千。如果真的出现这种情况，恐怕市场中的所有角色都将成为受害者。索罗斯在《金融炼金术》中对于银行不理性的放贷行为阐述道，只要他的业务量达到一定程度，政府就不会允许它宣布破产。如果是正面，他就赢了；如果是反面，他也不会输。在市场整体的不理性中，参与者几乎别无选择：要么随波逐流，要么退出竞争。如果管理层执意采取纠正措施，就要陷入一场突如其来的崩溃的风险。在美国2008年的次贷危机中，政府的表现似乎已经说明了一切。相对于“惩罚”而言，“救助”才是对所有人更有利的选择。恐怕这就是，即使第一次救市政策没有通过，第二次救市政策也必须通过的原因。发现没有？在参与者偏向形成下的价格，的确可以改变真实的价值！即使管理层投放更多的货币，也不会有魄力彻底否定集体的偏向。当然，市场中所有的角色，仍旧会在相当长时间内替这种偏向买单。同时，信贷扩张越是疯狂，市场就要用越长时间的萧条来消化。日本经济长期的萧条，或许就是真实的写照。

在读《金融炼金术》的时候，我始终猜测索罗斯对于“价值”的看法应该是多维的，否则很多问题都会变得难以理解。比如，抵押信托公司那种“庞氏骗局”式的规则，的确可以依赖不断涌入的后来者推高价格。高收益可以吸引更多接受溢价的买入者，而更多的溢价买入者又能够增加收益。我们谁也无法否认，抵押信托公司的收益的确是真实的，我们可以把这种高收益视为一种“微观价值”。但“庞氏骗局”的本质就注定了泡沫无法永远延续下去，最终走向破产的崩溃似乎又是在受到一种“宏观价值”的吸引。

我相信，很多人一定会认为我为“庞氏骗局”引入一个所谓的“微观价值”是没有意义的。毕竟，“庞氏骗局”本来就是完全没有价值的搏傻游戏。但事实上，任何经济领域的泡沫，都可以理解为是一种“庞氏骗局”的变种。难道信贷扩张与获利效应不是一种需要相互依赖的“庞氏骗局”吗？我们以房地产市场为例，难道生产钢筋、水泥的公司所实现的利润不是真实的吗？问题在于，这种真实的利润的确依赖于信贷扩张与获利效应的自我强化。房屋作为抵押品的价值与银行的放贷行为存在明显的反身性，这在《金融炼金术》中早有阐述。即使是贷款者对未来收入的预期，都会在很大程度上受到信贷扩张的影响。如果不是信贷创造出更多的货币，那么市场中任何角色赚钱都会变得困难得多。很显然，市场参与者（银行、上市公司、投资者、贷款者）所关注的报表上的所谓价值，大部分都属于“微观价值”。温家宝总理曾经说过：“通货膨胀就像一只老虎，如果放出来就很难再关进去。”由此可见，通货膨胀的确存在自我强化的效应，我们似乎看不到物价水平会受到所谓价值的制约。我们在报表中看到所谓的价值，完全都是以货币作为计价单位的。而货币本身，早已参与到这种自我强化的效应中来。

在某种自我强化的潮流中，我们似乎找不到任何可以制约这种潮流的因素。但当这种潮流最终完结以后，似乎价格又会受到“宏观价值”的吸引。我们很难相信，日本长达二十年的经济萧条，会与此前疯狂的经济泡沫无关。全部都是世界上史无前例的：高储蓄率、经济泡沫、经济萧条，先后发生在日本这个国家。很难让我们相信，储蓄、泡沫、萧条之间不存在相关性。虽然日本二十年前的经济泡沫能够在自我强化中几乎不受“宏观价值”的影响而走向疯狂，但疯狂过后仍旧无法摆脱大幅提前透支“宏观价值”而形成的萧条。或许，我们可以用“每股收益”之类的数据来形象地表述“微观价值”，那么我们又该如何表述“宏观价值”呢？对于这个问题，我只能说感觉它是存在的，但却无法将其精确地表述出来。最多，我只能在引入这个概念以后，讨论一下市场的特征罢了。有意思的是，虽然我无法对“宏观价值”进行表述，但我却可以简单地对其特点进行解释。事实上，我所谓的“宏观价值”，很类似于传统经济理论中的价值，或者说是大众对价值正统的理解方式。这不由让我想起了前面讲述的，那个美国黄豆市场数十年的走势图：微观上的追涨杀跌，可以造就宏观上的高抛低吸。或许对于“价值”也是同样的道理，微观上的“谬论”，放大到宏观上就是“真理”。

按照传统经济学理论的说法，大量的购买行为将抑制其后的购买行为，因为更多的购买行为可以推高价格。而按照需求曲线的说法，更高的价格会抑制需求，从而形成某种均衡。但在真实的市场中，我们的确无法看到这种“负反馈”的制约。大量的购买行为即使形成了更高的价格，仍旧会创造更多的购买行为。通货膨胀的现象并不能抑制通货膨胀的继续发展（负反馈），反而通胀却能够产生自我强化的效应（正反馈）。我相信，上述的说法在微观上是成立的。但在宏观上，通货膨胀的本身，就会为其后的紧缩埋下伏笔。或者说本轮通胀持续的时间与幅度越大，我们这些同龄人成为“紧缩的一代”的可能性就越大。我的结论是这样的：某种经济现象在小尺度上具有自我强化的正反馈效应；而在大尺度上具有自我抑制的负反馈效应。这或许就完美地解释了市场为什么总是不稳定的，也解释了趋势为什么是市场中的常态。按照传统经济学理论，如果一种经济现象的本身就能够（在微观上）抑制这种经济现象，那么趋势就应该是非常罕见的！其实，探讨某种经济现象究竟是“自我强化”，还是“自我抑制”的问题。是否类似于探讨光究竟是粒子还是波呢？波粒二象性告诉我们，站在不同的角度下，两个答案都是正确的。

我强烈建议大家阅读一下《金融炼金术》，以及一本叫做《临界》的书。我觉得，在本书中过多地引用上述两本书的观点，不如读者朋友自己亲自来阅读一下这两本不错的作品。前面所说的大尺度上表现出的自我抑制的负反馈效应，很类似于《临界》中讲述的封存效应。对于市场而言，正反馈是即刻兑现的，但负反馈却是被封存起来的。在一次壮观的趋势潮流中，市场是不会考虑累积起来的负反馈效应的，在这种情况下你考虑这种效应也只会被市场所嘲笑。只要音乐还没有停止，每位市场参与者就必须继续跳舞。封存起来的负反馈何时被释放出来，很类似于地壳运动被封存起来的能量何时引发地震一样。事实上，预测市场会发生什么事件，远比预测事件何时发生简单得多。我们可以合理地认为中国经济多年的扩张早晚都要出现“着陆”，但预测着陆的时间几乎是不可能的事情。繁荣没有标准的尺度，最多我们只能认为越壮观的繁荣，所能出现的次数就越少。另外可知的一点就是，繁荣与其后的萧条在尺度上具有一定的相关性。我们甚至不可能知道最终的着陆究竟是“软”还是“硬”，但我们却可以知道越温和的萧条就需要运行越长的时间。

我们继续以房地产市场为例，现在房屋出租的收益率基本上在2％左右，而无论是贷款利率还是存款利率都要大大高于2％。很显然，房地产市场无疑已经存在泡沫了！任何买房的决策除了心理需求以外，恐怕都是建立在房价继续上涨的前提下。而正是因为这个假设而买房的行为，在不断地支撑、甚至推高房价。房价的上涨早已在很大程度上依赖于预期产生的“鞋襻效应”。房价泡沫破灭的风险的确在不断累积，但何时这种风险会被释放出来几乎是不可知的。这很像《临界》中讲述的“堆沙游戏”，沙堆的崩溃是不需要特定原因的，但沙堆何时崩溃又是不可知的。

这就是本节开头所说的，我的基本分析方法基本上无法用于实战的部分原因。事实上，我对于市场预测的“噪音尺度”，要大于趋势跟踪系统所能够捕捉利润的尺度。我的意思是，假设在我预期的崩溃前，市场最后一次冲锋都很可能使我获得不小的利润。当然，除了思考的乐趣以外，我的分析还是有一个明显的价值。趋势跟踪系统回答了多与空的问题，而我的基本分析则是猜测那些市场更加不可能稳定。如果市场的价格完全是由于预期所支撑的话，那么虽然我无法回答继续上涨还是下跌的问题，但我却可以合理地认为这种市场处于非涨即跌的不稳定状态。我写本章的另一个目的则是告诉大家：对市场进行基本分析的时候，不要只考虑业绩之类的“微观价值”，而要对市场参与者偏向进行更多的研究。而体现交易者偏向最好的指标就是价格，在忽略市场真实价格走势的情况下研究基本分析是没有意义的。如果说研究技术分析的真谛在于寻找大众的本能偏向，那么研究基本分析的真谛就在于寻找大众的思维偏向。

\clearpage
